\documentclass[UTF8,11pt,a4paper,fontset=windows]{ctexart}
\usepackage[a4paper,left=2.2cm,right=2.2cm,top=2.2cm,bottom=2.2cm]{geometry}
\usepackage{booktabs}
\usepackage{fancyvrb}
\usepackage{fvextra}
\usepackage{hyperref}
\usepackage{xcolor}
\hypersetup{colorlinks=true,linkcolor=blue,urlcolor=blue}
\setmainfont{Times New Roman}
\setCJKmainfont{SimSun}[BoldFont=SimHei]
\setlength{\parindent}{2em}
\setlength{\parskip}{0.2em}
\VerbatimFootnotes

\title{\textbf{第17组\quad 完整原始数据}\\[0.4em]
\large 参数化 DCP 流式去雾硬件与精度--面积优化}
\author{许彦哲\quad 202618023429017\\ 中国科学院大学\ 前沿交叉科学学院}
\date{2026 年 11 月}

\begin{document}
\maketitle
\thispagestyle{empty}

\section{数据来源与工具}

本文件汇总 project2 的完整原始输出。所有质量、面积、真实图像指标均由 project2 工具重新运行得到；project1 的旧输出不作为依据。

\begin{center}
\small
\begin{tabular}{ll}
\toprule
项目 & 版本/路径\\
\midrule
Verilator & vUNKNOWN-built20261101\\
g++ & MinGW-W64 14.2.0\\
Yosys & 0.69 (git sha1 9f75ca1f9)\\
ABC & UC Berkeley ABC 1.01 (2026-11-09)\\
数据目录 & \texttt{project2/data}\\
证据目录 & \texttt{project2/docs/evidence}\\
\bottomrule
\end{tabular}
\end{center}

\section{文件清单与 SHA256}

下表中的 SHA256 覆盖输入、golden、最终输出与 RTL 源文件；其中 C++ 888/5 输出与 RTL 888/5 输出的哈希相同，说明二者逐字节一致。

\VerbatimInput[fontsize=\scriptsize,frame=single,breaklines=true,breakanywhere=true]{
../../project2/docs/evidence/sha256.txt}

\section{位宽配置的质量指标}

输入为 \texttt{data/real\_dcp/in.raw}（两帧相同图像），参考为 \texttt{data/dcp\_ref/golden\_video.raw}（888）。取第 1 帧输出计算 PSNR、SSIM、CIEDE2000 与 $Q$。原始输出如下。

\VerbatimInput[fontsize=\scriptsize,frame=single,breaklines=true,breakanywhere=true]{
../../project2/docs/evidence/quality_points.log}

\section{面积综合原始统计}

面积流程为 Yosys$\to$BLIF$\to$ABC，行缓存 \texttt{fifo\_line} 为 blackbox。ABC 的 \texttt{lat} 为寄存器数，\texttt{nd} 为 LUT6 数。四个配置的原始统计如下。

\VerbatimInput[fontsize=\scriptsize,frame=single,breaklines=true,breakanywhere=true]{
../../project2/docs/evidence/area_r8g8b8y8.log}

\VerbatimInput[fontsize=\scriptsize,frame=single,breaklines=true,breakanywhere=true]{
../../project2/docs/evidence/area_r8g8b8y7.log}

\VerbatimInput[fontsize=\scriptsize,frame=single,breaklines=true,breakanywhere=true]{
../../project2/docs/evidence/area_r8g8b8y6.log}

\VerbatimInput[fontsize=\scriptsize,frame=single,breaklines=true,breakanywhere=true]{
../../project2/docs/evidence/area_r8g8b8y5.log}

\section{实验 A 的完整最速上升日志}

设置 A：$Q>80$、面积权重 200。原始日志 \texttt{hill\_A.log} 如下。

\VerbatimInput[fontsize=\scriptsize,frame=single,breaklines=true,breakanywhere=true]{
../../project2/docs/evidence/hill_A.log}

\section{实验 B 的完整最速上升日志}

设置 B：$Q>90$、面积权重 120。原始日志 \texttt{hill\_B.log} 如下。其中 \texttt{lut=None} 表示该点被质量门限淘汰。

\VerbatimInput[fontsize=\scriptsize,frame=single,breaklines=true,breakanywhere=true]{
../../project2/docs/evidence/hill_B.log}

\section{真实图像指标}

对 \texttt{clear.raw} 计算 PSNR/SSIM。原始输出如下。

\VerbatimInput[fontsize=\scriptsize,frame=single,breaklines=true,breakanywhere=true]{
../../project2/docs/evidence/real_image_metrics.log}

\section{复现命令}

\begin{verbatim}
cd project2
python tools\dcp_opt.py --start 8,8,8,8 --workers 8 ^
  --qmin 80 --area-weight 200
python tools\dcp_opt.py --start 8,8,8,8 --workers 8 ^
  --qmin 90 --area-weight 120
\end{verbatim}

\end{document}
